% Figure --- the revocation exposure regime: which bound binds a stale card.
% Provenance: cor:fh-window (exposure <= T_p + T_g under a partition bound and
% a delivery service level) and thm:fh-conv item 3 (no bound without them);
% the card's TTL is the fail-closed ceiling. Drawn by hand; TTL = 8 rounds.
%
% Reader question: which bound actually stops a stale card from being
%   admitted: the operational bounds, or the card's own expiry?
% Claim: below the TTL diagonal the sum of the two operational bounds binds
%   exposure; above it those bounds exceed the TTL and only expiry binds.
% Evidence: Corollary~cor:fh-window (T_p+T_g) and Theorem thm:fh-conv item 3
%   (no bound without a partition bound), TTL=8 rounds (illustrative).
% Must distinguish: the region where the operational bound is tighter than
%   the TTL from the region where the TTL alone is doing the work.
% Grammar chosen: regime plane in rounds (tpl-quadrant.tex idiom, adapted to
%   a diagonal split rather than four cells, since the boundary is a sum, not
%   two independent yes/no axes). Grammar rejected: a timeline would need a
%   third axis for the partition bound; a bare formula restates the corollary
%   without showing which regime a deployment actually lands in.
% Five-second test: at partition bound 6 and delivery level 6, which bound is
%   doing the work?
\pdfigurehue{pdgold}%
\begin{figure}[H]
\centering
\begin{tikzpicture}[pd figure,x=0.76cm,y=0.40cm]
  \node[pd title,anchor=west,text width=9.6cm,align=left] at (0,12.2)
    {which bound binds a stale card: partition bound against delivery level};
  \draw[pd focus rule] (0,0) -- (8,0);
  \draw[pd focus rule] (0,0) -- (0,8);
  \draw[pd breach rule] (10,0) -- (10,10) -- (0,10) -- (0,8);
  \draw[pd focus rule,line width=1.4pt] (8,0) -- (0,8);

  \draw[pd rule,-{Stealth[length=1.8mm,width=1.1mm]}] (0,0) -- (10.4,0);
  \draw[pd rule,-{Stealth[length=1.8mm,width=1.1mm]}] (0,0) -- (0,10.6);
  \foreach \t in {2,4,6,8,10}
    {\draw[pd tick] (\t,0) -- (\t,-0.30);
     \node[pd label,anchor=north] at (\t,-0.42) {\t};
     \draw[pd tick] (0,\t) -- (-0.14,\t);
     \node[pd label,anchor=east] at (-0.22,\t) {\t};}
  \node[pd label,anchor=north,text width=6.5cm,align=center] at (5,-1.15)
    {partition bound, in rounds};
  \node[pd label,anchor=south,rotate=90,text width=6.5cm,align=center] at (-1.75,5)
    {delivery service level, in rounds};

  \node[pd focus label,anchor=south west] at (0.30,2.55)
    {exposure $\leq T_p+T_g$};
  \node[pd breach label,anchor=north west] at (1.60,9.55)
    {only the TTL (expiry) binds};
  \node[pd focus label,anchor=south west] at (4.55,5.10)
    {TTL: 8 rounds here};
\end{tikzpicture}
\caption{Which bound binds a stale card after revocation. Under a partition
bound and a delivery service level, exposure is at most the two added
together (Corollary~\ref{cor:fh-window}, below the gold TTL line); once that
sum exceeds the card's time to live the expiry binds instead and the harbor
fails closed (above the line, red dashed border). Without a partition bound the horizontal axis is unbounded and
only the expiry remains (Theorem~\ref{thm:fh-conv}, item 3); the line is
drawn for an illustrative time to live of eight rounds. [internal]}
\label{fig:fh-revocation-regime}
\end{figure}
